\section{הרחבה עתידית: \tech{Regime-aware Reinforcement Learning}}

לאחר שהממצאים מצביעים על כך שה־HMM שימושי יותר לתיאור מצב השוק מאשר לחיזוי נקודתי של היום הבא, עולה כיוון המשך טבעי: להשתמש במצב שה־HMM מעריך כחלק ממערכת שמקבלת החלטות לאורך זמן.

חשוב להבדיל בין שתי בעיות. ב־\tech{Supervised Learning} המודל מקבל תכונות ומנסה לנבא יעד ידוע, למשל האם היום הבא יהיה חיובי. ב־\tech{Reinforcement Learning (RL)} אין רק שאלת חיזוי. סוכן בוחר פעולה בכל שלב, מקבל תגמול, ולומד מדיניות שמטרתה לשפר את התוצאה המצטברת לאורך רצף של החלטות.

הרחבה זו לא בוצעה בפועל במסגרת הדוח. היא מוצגת כהצעת המשך שמבוססת על התוצאה המרכזית של הניסוי הקיים.

\subsection{ניסוח הבעיה כ־\tech{Markov Decision Process}}

ב־RL מקובל לתאר את הבעיה באמצעות \tech{Markov Decision Process (MDP)}. בכל יום הסוכן רואה מצב קלט, בוחר פעולה ומקבל תגמול.

אפשר להגדיר זאת כך:

\begin{itemize}
    \item \textbf{תצפית - \tech{Observation}.} ארבעת המאפיינים הקיימים, הסתברויות המצבים של ה־HMM, אנטרופיית ה־\tech{posterior}, והפוזיציה הנוכחית בתיק.
    \item \textbf{פעולה - \tech{Action}.} בגרסה פשוטה, \tech{sell / hold / buy}. בגרסה רציפה יותר, בחירת משקל \(w_t\in[0,1]\) שמייצג את שיעור החשיפה לנכס.
    \item \textbf{תגמול - \tech{Reward}.} תשואת התיק לאחר עלויות עסקה, עם אפשרות לקנס נוסף על סיכון או על תחלופה גבוהה של הפוזיציה.
\end{itemize}

ניסוח אפשרי לתגמול הוא:

\begin{equation}
R_{t+1}=w_t r_{t+1}-c\lvert w_t-w_{t-1}\rvert-\lambda\,\mathrm{RiskPenalty}_{t+1}.
\end{equation}

האיבר הראשון הוא תשואת הפוזיציה. האיבר השני מקזז עלות עסקה כאשר משנים את החשיפה, כאשר \(c\) מייצג את גובה העלות. האיבר השלישי מאפשר להעניש התנהגות מסוכנת, בהתאם למשקל \(\lambda\). בניסוי כזה עלויות עסקה אינן פרט טכני: מדיניות שמבצעת הרבה פעולות יכולה להיראות טובה לפני עלויות ולהפוך לחלשה לאחריהן.

\subsection{מדוע לשלב HMM בתוך RL?}

הרעיון אינו להחליף את ה־HMM ב־RL, אלא לתת לסוכן את אומדן מצב השוק שה־HMM כבר למד. במקום לספק רק תווית קשיחה כמו ``מצב 3'', אפשר להזין את וקטור ההסתברויות המלא ואת רמת אי־הוודאות:

\begin{equation}
\left[\gamma_t(1),\ldots,\gamma_t(K),H_t\right].
\end{equation}

כך הסוכן מקבל שני סוגי מידע: איזה משטר נראה סביר כרגע, ועד כמה ה־HMM בטוח בכך. לדוגמה, ייתכן שמדיניות טובה תפעל אחרת כאשר מצב אחד מקבל הסתברות של \pct{98}, לעומת יום שבו שני מצבים מקבלים הסתברויות דומות.

יתרון חשוב של גישה זו הוא שאין צורך לקודד ידנית כלל כמו ``במצב משברי מוכרים''. הסוכן מקבל את המידע ולומד מתוך פונקציית התגמול האם וכיצד להשתמש בו.

הניסוי הנקי ביותר יהיה ניסוי \tech{ablation}, כלומר השוואה שבה מוסיפים רכיבי מידע בהדרגה:

\begin{enumerate}
    \item \textbf{\tech{RL-Features}.} הסוכן רואה רק את ארבעת המאפיינים המקוריים.
    \item \textbf{\tech{RL-Regime}.} הסוכן רואה את אותם מאפיינים ובנוסף את הסתברויות המצבים של ה־HMM.
    \item \textbf{\tech{RL-Regime-Uncertainty}.} נוסף גם מידע על אנטרופיית ה־\tech{posterior} ועל הפוזיציה הנוכחית.
\end{enumerate}

כך אפשר לענות על שאלה ממוקדת: האם המידע החבוי על משטר השוק מוסיף ערך מעבר למאפיינים הגולמיים שכבר קיימים.

\subsection{אלגוריתמים אפשריים}

לפעולות בדידות אפשר להתחיל במודל \tech{actor-critic} פשוט. אם הפעולה היא משקל רציף בתיק, שתי אפשרויות טבעיות הן \tech{PPO} \cite{schulman2017ppo} ו־\tech{Soft Actor-Critic (SAC)} \cite{haarnoja2018sac}.

\tech{PPO} נפוץ משום שהוא מגביל את גודל עדכוני המדיניות ומנסה לשמור על אימון יציב יחסית. \tech{SAC} מתאים במיוחד לפעולות רציפות ומשתמש באנטרופיה כדי לעודד חקירה. קיימות גם מסגרות כמו \tech{FinRL}, שמספקות תשתית מוכנה יחסית לנתונים, סביבות מסחר ועלויות שוק \cite{liu2021finrl}. עם זאת, ספרייה מוכנה אינה פותרת את בעיות המתודולוגיה: עדיין צריך להגדיר חלוקה כרונולוגית, עלויות, מדדי הערכה ובקרת זליגה.

\subsection{מה מראה הספרות?}

הספרות על RL למסחר מציגה תוצאות מעורבות. Yang et al. אימנו שילוב של \tech{PPO}, \tech{A2C} ו־\tech{DDPG} על 30 מניות במדד Dow Jones ודיווחו על תוצאות טובות יותר בחלק ממדדי הסיכון והתשואה בניסוי שלהם \cite{yang2020ensemble}. Zhang, Zohren and Roberts בחנו \tech{Deep RL} על 50 חוזים עתידיים ממספר סוגי נכסים ודיווחו על ביצועים חיוביים גם כאשר נכללו עלויות עסקה \cite{zhang2019drltrading}.

מצד שני, תוצאות חיוביות אינן מובטחות כאשר הפרוטוקול מחמיר. Kruthof and M{\"u}ller בחנו \tech{SAC} על שבעה מאגרי מניות וביותר מ־300 שנות \tech{out-of-sample} מצטברות. הם לא מצאו יתרון שיטתי על פני כלל \(1/N\), ועלויות עסקה מתונות פגעו במיוחד במדיניות בעלת תחלופה גבוהה \cite{kruthof2025beat1n}.

לכן המסקנה מהספרות אינה ש־RL בהכרח עובד או בהכרח נכשל. היא מדגישה שהערכת RL פיננסי רגישה מאוד להגדרת הניסוי, לעלויות, לתקופת הבדיקה ולבחירת מודלי הייחוס.

\subsection{כיצד לבצע ניסוי המשך הוגן?}

אם מרחיבים את הפרויקט בעתיד, כדאי לשמור על אותם עקרונות שהנחו את הניסוי הנוכחי:

\begin{itemize}
    \item חלוקה כרונולוגית ו־\tech{walk-forward}, ללא ערבוב אקראי.
    \item התאמת ה־HMM וה־\tech{scaler} רק על המידע הזמין בכל חלון זמן.
    \item מספר אתחולים לכל מדיניות, משום שאימון RL רגיש לאקראיות.
    \item עלויות עסקה ו־\tech{turnover} מפורשים בתגמול ובהערכה.
    \item השוואה ל־\tech{Buy-and-Hold}, ממוצע נע, מדיניות ללא HMM ומדיניות עם HMM.
    \item מדדים כגון תשואה מצטברת, \tech{Sharpe ratio}, \tech{maximum drawdown} ו־\tech{turnover}, לצד בדיקה אם התוצאה נשמרת לאחר עלויות.
\end{itemize}

לא בוצע ניסוי RL בדוח זה, ולכן אין כאן טענה שמידע על משטרים ישפר בפועל את ביצועי המסחר. ההצעה היא שאלה מחקרית להמשך: אם ה־HMM אכן טוב יותר בזיהוי מצב מאשר בחיזוי יום יחיד, כדאי לבדוק האם אומדן המצב הזה מועיל כאשר הבעיה מוגדרת כקבלת החלטות סדרתית.
